\section{Discussion}

\method{} does not prove that deleted data has no possible remaining effect. Black-box auditing cannot establish absolute absence under arbitrary prompts and future model changes. Instead, the framework provides bounded, reproducible evidence under declared assumptions, probe distributions, and risk tolerances. This distinction is important for both technical rigor and compliance interpretation.

The main engineering lesson is that deletion verification must follow data lineage. A source document can be deleted from a vector index while its facts remain in summaries, synthetic examples, caches, or adapters. Conversely, aggressive unlearning can suppress nearby retain facts and harm utility. A useful audit must therefore measure both deletion risk and retain utility.

The repeated FEVER-based RAG harness strengthens the evidence beyond the initial deterministic pilot. The advanced experiments further add low-rank adapter training, PEFT model wrapping, a three-seed pretrained SmolLM2 LoRA causal-LM run, a larger Qwen2.5-0.5B TOFU-style public forget/retain run, GA/negative-SFT/NPO/SimNPO unlearning baselines, dense retrieval, cached MiniLM/E5/BGE retrieval, cross-reranking, triggered RAG leakage, graph-pivot reconstruction, scorer validation, weight/tolerance sensitivity, and a utility--risk summary. Two limitations remain. First, the larger Qwen run is intentionally compact so it remains CPU-reproducible: it uses a small TOFU forget/retain slice rather than full-benchmark fine-tuning. Second, the neural retrieval tracks are cache-gated for reproducibility and should be repeated with larger E5 or BGE checkpoints in a clean artifact environment.

\textbf{Threats to validity.}
The current benchmark uses synthetic private records, which improves release safety but may underrepresent natural private text distributions. The scorer validation table is therefore a labeled synthetic diagnostic, complementing the human adjudication on natural private text; the artifact now includes a human-adjudication log for redacted deployment samples. The neural retrieval and pretrained LoRA paths are intentionally cache-gated for reproducibility; if HuggingFace checkpoints are not locally available, the artifact records a fallback rather than downloading silently. Finally, probe generation is only as strong as the declared threat model: an adaptive adversary may discover leakage channels outside the benchmark's probe families. Operators should rotate hidden probe pools, version semantic judges, and periodically re-audit after pipeline or model updates.

\textbf{Pure black-box operation.}
When retrieval metadata is unavailable, \method{} can still audit answer-side direct, paraphrase, counterfactual, extraction, and watermark evidence. The answer-only experiment intentionally removes artifact identifiers and citations, so retrieval dependence is no longer credited unless it manifests in the text. This setting is weaker than gray-box RAG auditing, but it avoids assuming privileged access and gives practitioners a conservative deployment mode.

\textbf{Formal guarantees.}
ADCU is not a replacement for certified removal, deletion-as-control, differential privacy, or pan-private continual release. Its role is complementary: it audits deployed pipeline behavior, derivative artifacts, caches, summaries, adapters, and retrieval-control layers that often sit outside the formal guarantee boundary. When a component has a formal deletion guarantee, ADCU should treat that guarantee as prior evidence and spend more probes on uncovered downstream artifacts.
